\chapter{Analysis and Verification}
In this chapter we will use the methodology developed in \Cref{ch:1} to evaluate model-based designs of embedded systems.
For this, we first have to define what to evaluate for--what makes a model \emph{correct}.

\section{Correctness and Model Checking}

\begin{definition}[Correctness, Specification, Validation, and Verification]
    We define a model to be \emph{correct} if it meets the right \emph{specifications} in its operating environment.

    The specification of a model is a mathematical statement of the design objectives---the desired properties of the system.

    To \emph{validate} a model is to prove if its specification meets the environments needs in the first place---\emph{you build the right thing}.

    Finally, the \emph{verification} of a model is to check if the system has been build according to its specifications---\emph{you build the thing right}.
\end{definition}

As previously hinted, the goal of a model-based design approach is to not only synthesise software for embedded systems from a model, but also automatically verify it in doing so.
The latter can be done with \emph{model checking}, a process which determines whether a system satisfies a formal specification.
The specification is a \emph{temporal logic} formula, which expresses the desired systems properties over time.
To define such formulas, we have to formally define \emph{execution traces} and \emph{propositional logic formulas}.

\begin{definition}[Execution Traces and Propositional Logic]
    An execution trace of an SM is a sequence of the form \(q_0, q_1, q_2, \dots\) with \(q_i = (x_i, s_i, y_i)\) with \(x_i\) being the input valuation, \(s_i\) the state, and \(y_i\) the output valuation at step \(i\).

    A propositional logic formula is an extended version of the previously defined guards with \(s\) being another atomic statement, meaning that the SM is in state \(s\).
    A more complicated formula is e.g., \(s_0 \wedge x = 1 \wedge \neg y\), meaning that \(s_0\) is the active state, the input \(x\) is present with value \(1\) and the output is not present during the reaction.

    We define a propositional logic formula to \emph{hold} for a trace \(q_0, \dots\) iff it holds for \(q_0\).
\end{definition}

The reason in the last definition is, that we can define adequate propositional logic formulas to reason, that the SM must be correct starting from the first step of the trace.

There are two main variants of temporal logic formulas we will go into, linear temporal logic (LTL), and real-time temporal logic.

\newcommand*{\G}{\textbf{G}}
\newcommand*{\F}{\textbf{F}}
\newcommand*{\X}{\textbf{X}}
\newcommand*{\U}{\textbf{U}}
\subsection{Linear Temporal Logic}
LTL formulas are an extension of propositional logic formulas with the following grammar.
\begin{definition}[Linear Temporal Logic Formula]
    An LTL formula extends propositions with four operators.
    The simplest LTL formula is just a proposition \(p\).

    For LTL formulas \(\phi, \phi_1, \phi_2\), we define the following LTL formulas for a trace \(q_0, q_1, \dots\).
    \begin{IEEEeqnarray*}{l'l}
        \G\phi & \text{\(\phi\) holds for every suffix of \(q_0, q_1, \dots\).} \\
        \F\phi & \text{\(\phi\) holds for some suffix of \(q_0, q_1, \dots\).} \\
        \X\phi & \text{\(\phi\) holds for \(q_1, \dots\).} \\
        \phi_1\U\phi_2 & \text{\(\phi_1\) holds for all suffixes of \(q_0, q_1, \dots\) until a suffix for which \(\phi_2\) holds.}
    \end{IEEEeqnarray*}
\end{definition}

Thus, \(\G\phi\) is equivalent to mean, that \(\phi\) holds for all steps of the trace, while \(\F\phi\) only needs one step to hold to fulfil the formula.
\(\G\F\phi\) means, that \(\phi\) holds infinitely often.
\(\F\G\phi\) means, that eventually \(\phi\) holds for every step.
\(\G(\phi \Rightarrow \F \psi)\) means, that after \(\phi\) holds, \(\psi\) must hold eventually.
\(\G(\phi \Rightarrow \X\X\psi)\) means, that after \(\phi\) holds, \(\psi\) holds exactly two reactions after.

\begin{lemma}[Temporal Operators and Relationships]
    For some trace and LTL formula \(\phi\), we have
    \begin{IEEEeqnarray}{l}
        \G\phi = \neg\F\neg\phi, \\
        \F\phi = \true \U \phi, \\
        \text{\(X\) can not be expressed in terms of \(\G\), \(\F\), or \(\U\).}
    \end{IEEEeqnarray}
\end{lemma}

Real-time temporal logic redefines the syntax of \(\F\) to uphold deadlines, e.g., \(\F_{\le 2.5 \text{ms}}\) means, that the following formula must hold within \(2.5\) ms of the trace.

\section{Equivalence and Refinement}
Another analysis goal for this chapter is the comparison of different models.
There are different equality metrics, which can be used to compare two models.

\begin{definition}[Type Equivalence]
    Two actors are \emph{type equivalent} if they have the same input and output ports, which have the same sets of possible input and output valuations.
    In other words, they typecheck.
\end{definition}

\begin{definition}[Language Equivalence]
    Two actors are \emph{language equivalent} if for every input sequence they can produce the same output sequence.
\end{definition}

Unequivalent machines can still be \emph{refinements} of each other, which is understood as a onesided equivalence.
E.g., a machine can be a \emph{type}, \emph{language}, or \emph{simulation refinement} of another machine, if it is type compatible, produces only outputs which the other can produce given the same input sequence, or produces the same output values on each reaction.

To formally define the refinements, recall SMs \(A\) and \(B\) with with states \(S_A, S_B\), initial states \(I_A, I_B\), input/output valuation sets \(V_A, V_B\) for set of input signals \(P_A, P_B\), set of output signals \(Q_A, Q_B\), and update functions \(u_A, u_B\).

\begin{definition}[Type Refinement]
    \(B\) is a type refinement of \(A\) if
    \begin{IEEEeqnarray}{ll"l}
        & P_B \subseteq P_A & \text{\(B\) does not need all input signals,} \\
        & Q_A \subseteq Q_B & \text{\(B\) may have more output signals,} \\
        \forall p \in P_B: & V_{Ap} \subseteq V_{Bp} & \text{input valuations typecheck,} \\
        \forall q \in Q_B: & V_{Bq} \subseteq V_{Aq} & \text{output valuations typecheck.}
    \end{IEEEeqnarray}
    In this case, \(B\) can replace \(A\) without causing a type conflict.
\end{definition}

\begin{definition}[Behaviour and Language Refinement]
    The \emph{behaviour} of a SM is a tuple of input signals and produced output signals, such that the output signal matches the output sequence generated by an input sequence.
    The \emph{language} set \(L(\cdot)\) of a SM is the set of all behaviours.

    \(B\) is a language refinement of \(A\) if \(L(B) \subseteq L(A)\).
    I.e., \(B\) can replace \(A\) without producing anything \(A\) could not have produced.
\end{definition}

\begin{definition}[Simulation Relation / Simulation Refinement]
    \(A\) \emph{simulates} \(B\), if there is a relation \(S \subseteq S_B \times S_A\) which
    \begin{IEEEeqnarray}{c}
        (I_B, I_A) \in S \\
        \forall (s_B, s_A) \in S, i \in V^{in}_B, (s_B', o) \in u_B(s_B, i): \exists (s_A', o) \in u_A(s_A, i): (s_B', s_A') \in S
    \end{IEEEeqnarray}
    In other words, \(A\) simulates \(B\) if \(A\)s initial state simulates \(B\)s, and for each possible input \(B\) can have, the resulting state is also simulated by \(A\) and has the same output.
    \(S\) serves as a composed state machine of \(A\) and \(B\) that exists exactly when \(A\) simulates \(B\).

    Because \(A\) must only accept inputs which \(B\) accepts, we can show that \(L(B) \subseteq L(A)\) in this case.
    Thus, if \(A\) simulates \(B\), \(B\) is a simulation refinement of \(A\) and \(B\) is a language refinement of \(A\).
\end{definition}

\begin{definition}[Bisimulation]
    An actor \(M_2\) is \emph{bisimilar} to actor \(M_1\) if for every state of \(M_1\) there exists a state in \(M_2\) which reacts exactly the same way.

    For deterministic SMs, language equivalence and bisimilarity are the same, but this is not the case in nondeterministic SMs.

    The formal definition resembles the simulation relation, but with the difference, that the simulation must simulate both SMs at the same time.
    This is harder to that \(A\) and \(B\) simulate each other, as the simulation relations may not equal, which is forced in this definition.
\end{definition}

\section{Reachability Analysis and Model Checking}
The goal in this section is to check for SMs if a target state \(s\) is reachable in the environment from initial state \(I\)---this is called an \emph{reachability problem}.

The approach to do this is to create SMs for the system under development (SUD) and environment, and composing them, such that input and output signals are connected.
This composed actor now has to be analysed.
In this case, the system is called \emph{closed} as the environment model provides all inputs and takes all outputs of the SUD SM.

Another analysis goal is to define an LTL formula for the composed actor, and automatically generate proofs/refutations.

Another task in the automated analysis is the search for simpler refinement actors to reduce state space and make other analyses easier.

\subsection{Graph Traversal Algorithms}
The problem with naive DFS and BFS is, that the runtimes and memory requirements are just too extreme for big state machines.
The advantage of BFS over DFS is, that it finds the quickest path to a failure, if it exists.

There is a nice variant of BFS which can be used to represent everything cleaner and with less memory.
If the graph to discover has nodes \([n]_0\), we can store a binary number \lstinline|M| \(\in \{1, 0\}^n\) which represents the current frontier in the search.
Checking if a node \(i\) was already discovered can be done with \lstinline|M && 2 << i|, and in this way, all neighbours of all new frontier nodes can be checked in one boolean operation.

\subsection{Using Temporal Logic}
As another approach, reachability analysis for state \(s\) can also be implemented using LTL formulas with \(\F s\), meaning, that for some suffix, i.e., eventually, \(s\) is reached.

LTL can also be used to specify the order in which some states are reached, e.g., \(\F(s_1 \wedge \F(s_2 \wedge \F(s_3 \dots)))\).

\subsection{Liveness Properties and Finding Acceptance Cycles}
As important as checking if bad conditions are not met, we can also check if positive conditions will be met.
E.g., \(\Phi := \F \G \phi\), after some time a desired \(\phi\) is always met.

Such a check can also be found by defining the \emph{system automaton} of \(\Phi\), constructing the synchronous \emph{Büchi automaton} composition of both SMs, and checking if there does not exists a reachable infinite cycle where the \(\phi\) state is reached.
In this case, we can refute \(\Phi\).
